%% ── 5-Р БҮЛЭГ: ДҮГНЭЛТ ─────────────────────────────────────
\chapter{Дүгнэлт}
\label{ch:conclusion}

\section{Судалгааны гол дүгнэлт}

Энэхүү судалгааны ажлын хүрээнд Flutter болон Node.js технологийг ашиглан
хэрэглэгчийн тэмдэглэлийн агуулга болон сэтгэлийн байдалд тохирсон ухаалаг
тэмдэглэлийн дэвтрийн мобайл аппликейшныг амжилттай боловсруулав.

Тодруулбал, дараах зорилтуудыг бүрэн хэрэгжүүллэв:

\begin{enumerate}
  \item Flutter (Dart) ашиглан Android болон iOS-д ажиллах мобайл аппликейшн бүтээв.
  \item Node.js + Express.js-ээр 6 REST API endpoint бүхий backend хөгжүүллэв.
  \item MongoDB Atlas үүлэн өгөгдлийн санд User болон Note өгөгдлийн загваруудыг хэрэгжүүллэв.
  \item JWT болон bcrypt ашиглан хэрэглэгчийн баталгаажуулалт, нууцлалыг хангав.
  \item flutter\_secure\_storage ашиглан токеныг аюулгүй хадгалах системийг нэвтрүүллэв.
  \item Mood tracking функцыг 5 төрлийн сэтгэл хөдлөлтэйгээр хэрэгжүүллэв.
  \item \textbf{AI sticker үүсгэлт} -- тэмдэглэлийн агуулгад тохирсон, овөрмөц стикер дизайн автомат үүсгэүлэв.
  \item \textbf{Calendar view} -- үүсгэсэн стикерийг календарийн тодорхой өдөрт холбож харуулах боломж олголов.
  \item \textbf{Push notification} -- тухайн өдрийн тэмдэглэлд суурилсан сануулга автоматаар илгээх систем нэвтрүүлэв.
  \item 15 хэрэглэгчээр туршиж 4.5/5-ийн сэтгэл ханамжийн үнэлгээ авав.
\end{enumerate}

\section{Шинэлэг байдал}

Манай системийн шинэлэг байдал нь:
\begin{itemize}  \item \textbf{AI суурилсан стикер} -- тэмдэглэлийн агуулга, mood дараах автоматаар стикер дизайн үүсгэдэг
  \item \textbf{Calendar view} -- үүсгэсэн стикерийг календарийн тодорхой өдөрт тавиж, визуалчлан харуулах
  \item \textbf{Ухаалаг push notification} -- календарийн темдэглэлд суурилсан push notification-ийг систем өөрөө автомат илгээдэг  \item \textbf{Mood-aware тэмдэглэл} -- хэрэглэгчийн сэтгэлийн байдлыг тэмдэглэлтэй холбосон
  \item \textbf{SplashRouter} -- токен-д суурилсан автомат нэвтрэх систем
  \item \textbf{Монгол хэрэглэгчдэд зориулсан} -- хэлний онцлог, хэрэгцээнд нийцсэн UX
  \item \textbf{Нээлттэй эхийн шийдэл} -- GitHub-т байрлуулсан, бусад хөгжүүлэгчид ашиглах боломжтой
\end{itemize}

\section{Хязгаарлалт болон цаашид хийх ажил}

\subsection{Одоогийн хязгаарлалтууд}

\begin{itemize}

  \item Офлайн горимд тэмдэглэл үүсгэх боломжгүй
  \item Зөвхөн нэг хэрэглэгчийн мэдээллийн сангийн бүтэц -- хамтарсан тэмдэглэл байхгүй
\end{itemize}

\subsection{Цаашид хийх ажлын чиглэл}

\begin{itemize}
  \item Тэмдэглэл хайх, шүүх, tag функц нэмэх
  \item iOS App Store болон Google Play дээр нийтлэх
  \item Admin dashboard хөгжүүлэх
  \item Офлайн горим (local SQLite/Hive)
\end{itemize}

\section{Судалгааны ажлын арга зүйн ач холбогдол}

\subsection{Технологийн сонголтын үндэслэл}

\textbf{Flutter} сонгосон шалтгаанууд:
\begin{itemize}
  \item Нэг кодын баазаас Android болон iOS-д ажиллах
  \item Google-ийн идэвхтэй дэмжлэг, том community
  \item Hot Reload функцоор хөгжүүлэлтийн хурд өндөр
  \item Material Design 3 дэмжлэг
\end{itemize}

\textbf{Node.js + Express.js} сонгосон шалтгаан:
\begin{itemize}
  \item JavaScript-ийн нэг хэлийг frontend болон backend-д ашиглах
  \item Async I/O-ийн улмаас олон хүсэлтийг зэрэг боловсруулж чадна
  \item npm-ийн баялаг экосистем (jwt, bcrypt, mongoose)
  \item RESTful API барихад хялбар
\end{itemize}

\textbf{MongoDB Atlas} сонгосон шалтгаан:
\begin{itemize}
  \item Schema-flexible -- тэмдэглэлийн агуулга тодорхой бус байж болно
  \item JSON-тэй ижил BSON бүтэц -- JavaScript объекттэй шууд нийцдэг
  \item Atlas -- үнэгүй tier байдаг, суулгалт хэрэгтэйгүй
  \item Геолокацийн дагуу server байршил сонгох боломж
\end{itemize}

\subsection{Хямдрал, уян хатан байдал}

Нээлттэй эхийн технологиудыг ашигласан тул:
\begin{itemize}
  \item Лиценз төлбөр байхгүй
  \item Эх кодыг бүрэн удирдаж болно
  \item Дараагийн хөгжүүлэлтэд чөлөөтэй нэмж өргөтгөх боломжтой
  \item GitHub-т нийтлэн бусад хөгжүүлэгчдэд нэмэлт хувь нэмэр оруулах боломж
\end{itemize}

\section{Хэрэглэгчийн туршилтын дэлгэрэнгүй дүн шинжилгээ}

15 оюутан хэрэглэгчийг гарын авлага, заавар байхгүй нөхцөлд апп-ыг ашиглуулж
дараах хэсгүүдэд үнэлгээ авав:

\subsection{Туршилтын дэлгэрэнгүй}

\textbf{Туршилтын хэсэг 1: Бүртгэл болон нэвтрэлт}
\begin{itemize}
  \item 15-аас 15 нь (100\%) нэг оролдлогоор амжилттай бүртгэж нэвтэрсэн
  \item Алдаатай имэйл оруулахад алдааны мэдэгдэл зөв гарсан
  \item Нэвтрэлтийн дундаж хугацаа: 2.3 секунд (сүлжээ хамгийн удаан нь)
\end{itemize}

\textbf{Туршилтын хэсэг 2: Тэмдэглэл үүсгэх}
\begin{itemize}
  \item 15-аас 15 нь (100\%) тэмдэглэл амжилттай үүсгэсэн
  \item Mood сонгох функцийг 14 нь (93\%) ашигласан
  \item Тэмдэглэл хадгалагдах хугацаа дундажаар 1.1 секунд
\end{itemize}

\textbf{Туршилтын хэсэг 3: Өдрийн тэмдэглэлийн жагсаалт}
\begin{itemize}
  \item Өнөөдрийн тэмдэглэлүүд зөв шүүгдэн гарч байсан
  \item Dashboard greeting цагийн дагуу зөв харагдсан
  \item 14-аас 15 хэрэглэгч (93\%) UI хялбар гэж үнэлсэн
\end{itemize}

\subsection{Туршилтын санал хүсэлт}

\begin{itemize}
  \item \textit{"Хялбар, гоё харагддаг апп байна"} -- Оюутан \#3
  \item \textit{"Mood feature сонирхолтой, ашиглахад хялбар"} -- Оюутан \#7
  \item \textit{"Дараагийн хувилбарт хайлт нэмчихвэл сайн байх"} -- Оюутан \#11
  \item \textit{"Dark theme нүдэнд тааламжтай"} -- Оюутан \#15
\end{itemize}
